% ============================================================
% 50_typography_semantics.tex — Schriftmakros für Box-Typen & semantische Inhaltsmakros
% ============================================================

% ── Die eine Rechenstelle für jede benannte Schriftrolle ─────────────
% Jede Grösse unten wird GERECHNET, nicht geschrieben: Sie ist ihr bp-Mass
% mal dem Rollenfaktor mal der Grundgrösse der ZSF (\ZSFBodyScale aus
% 30_layout_spacing, 1 solange \ZSFBodyFontSize nicht steht). Über \the und
% damit in pt, weil \fontsize sein Argument sonst als Zahl ohne Einheit lesen
% müsste — der Faktor kommt so als fertiges Mass an und nicht als Term.
%
% \ZSFBodyScale steht hier fest verdrahtet und nicht als vierter Parameter:
% Es gibt keine Schriftrolle im Textkörper, die der Grundgrösse NICHT folgen
% soll. Wäre sie ein Argument, wäre die Frage bei jeder neuen Rolle wieder
% offen — und genau eine Rolle bekäme sie irgendwann anders beantwortet.
\makeatletter
\newlength{\ZSF@fontSize}
\newlength{\ZSF@fontSkip}
\newcommand{\ZSF@scaledFont}[3]{%
  \setlength{\ZSF@fontSize}{\ZSFBodyScale\dimexpr #1\dimexpr #2\relax\relax}%
  \setlength{\ZSF@fontSkip}{\ZSFBodyScale\dimexpr #1\dimexpr #3\relax\relax}%
  \fontsize{\the\ZSF@fontSize}{\the\ZSF@fontSkip}\selectfont}

% Explizite bp-Grössen (wie Analysis) für präzise, font-unabhängige Kontrolle.
% Der Rollenfaktor ist 1, wo die Rolle keinen eigenen Regler hat — nur die
% Diagramm-Beschriftung bringt mit \ZSFDiagramLabelScale einen mit.
\newcommand{\ZSFfontChapter}{\sffamily\bfseries\ZSF@scaledFont{1}{8.1bp}{9.1bp}}
\newcommand{\ZSFfontSection}{\sffamily\bfseries\ZSF@scaledFont{1}{8.1bp}{9.1bp}}
\newcommand{\ZSFfontSubsubsection}{\sffamily\bfseries\ZSF@scaledFont{1}{7.2bp}{8.0bp}}
\newcommand{\ZSFfontBoxTitle}{\sffamily\bfseries\ZSF@scaledFont{1}{8.1bp}{9.1bp}}
% chktex 6
\newcommand{\ZSFfontNote}{\sffamily\itshape\ZSF@scaledFont{1}{6.7bp}{7.5bp}}
% Rechtsbuendiges Meta-/Kategorie-Tag im Titelbalken (vgl. \ZSFTitleTag in
% styles/60_boxes.tex). Etwas kleiner und kursiv, damit es den Titel begleitet
% statt mit ihm zu konkurrieren.
\newcommand{\ZSFfontTitleTag}{\sffamily\itshape\ZSF@scaledFont{1}{6.4bp}{7.0bp}}
\newcommand{\ZSFfontTitleHeader}{\sffamily\bfseries\ZSF@scaledFont{1}{11bp}{12.5bp}}
\makeatother

% ── Seitenmöbel: bewusst OHNE \ZSFBodyScale ──────────────────────────
% Fusszeile und Identity-Marker gehören zur SEITE und nicht zur Textspalte.
% Die Seitenzahl am Fuss ist eine feste physische Grösse — sie schrumpft nicht
% mit, nur weil der Text der ZSF enger gesetzt ist, und der 4pt-Marker wäre
% eine Stufe kleiner nicht mehr lesbar. Deshalb rohes \fontsize: Das ist hier
% die Aussage »absolut gemeint«, nicht ein vergessener Regler.
\newcommand{\ZSFfontFooter}{\sffamily\bfseries\fontsize{9.5pt}{9.5pt}\selectfont}
\newcommand{\ZSFfontIdentityMarker}{\fontsize{4pt}{4pt}\selectfont}

% Schrift/Farbe für Code-Kommentare — nur für das optionale Code-Modul
% (styles/67_code_comments.tex) benötigt. ce_green stammt aus
% styles/40_colors_structure.tex. Schadlos wenn das Modul deaktiviert ist.
\newcommand{\ZSFfontCodeComment}{\ttfamily\color{ce_green}}

% Semantische Inhalts-Fontgrössen — die zwei Stufen des Reglers `font`
% (`normal` | `dense`). Kapitel wählen sie nicht direkt, sondern über den
% Regler an der Box (styles/60_boxes.tex) oder an der ZSFtable
% (styles/20_tables.tex). `dense` ist die eine Stufe kleiner: lange Register,
% Vergleichsmatrizen, gedrängte Aufzählungen.
%
% EIN Paar für beide Bausteine, nicht je eines: »eine Stufe kleiner« ist
% dieselbe Entscheidung, ob sie eine Tabelle oder einen Boxinhalt trifft. Mit
% zwei Paaren müsste jede Änderung am einen im anderen nachgezogen werden
% (styles/README → Vereinigungstest).
%
% Beide holen die Display-Skips zurück: \normalsize und \footnotesize setzen
% sie auf die Klassenwerte, und die ZSF-Entscheidung dazu ist eine andere
% (\ZSFApplyDisplaySkips in 30_layout_spacing). Ohne diesen Nachsatz wären
% Formeln in jeder Box luftiger als im Fliesstext — lautlos, weil beides
% für sich korrekt aussieht.
% ── Der Rollenfaktor des Bausteininhalts ────────────────────────────
% Das Gegenstück zu \ZSFFlowTextScale, und aus demselben Grund: Wie laut steht
% der Inhalt der Bausteine — Boxtext, Tabellenzelle, Formel — neben den
% Titelbalken, die ihn gliedern?
%
% Ohne diesen Faktor gab es dafür nur \ZSFBodyFontSize, und das nimmt die
% Balken MIT. Der Zustand »Balken auf voller Grösse, Inhalt eine Stufe
% darunter« war damit nicht ausdrückbar — obwohl er der natürliche ist: Die
% Balkengrössen stehen absolut in bp (oben) und sind auf die Spaltenbreite
% abgestimmt, der Inhalt dagegen muss in eine Tabellenzelle passen. Das sind
% zwei Fragen, und sie hatten einen gemeinsamen Regler.
%
% Gemessen, woran das auffiel: Eine Formelzelle einer zweispaltigen Tabelle
% ist 86.2 pt breit. \sin(\arctan(x)) = x/\sqrt{1+x^2} braucht bei einem
% Inhalt von 8 pt 86.3 pt und bricht um; bei 7 pt sind es 75.9 pt. Wer das
% ueber \ZSFBodyFontSize loest, schrumpft die Balken von 8.1 auf 7.1 bp mit
% und aendert das Satzbild an einer Stelle, an der nichts falsch war.
%
% Angewandt NACH der Klassenstufe, nicht statt ihr: \normalsize und
% \footnotesize tragen \ZSFBodyScale bereits (30_layout_spacing rechnet die
% Stufen um, sobald \ZSFBodyFontSize steht). Der Rollenfaktor multipliziert
% nur noch — sonst zaehlte die Grundgroesse doppelt.
%
% Ein Faktor fuer beide Stufen, nicht zwei: »eine Stufe kleiner« (`font=dense`)
% ist eine Entscheidung PRO STELLE und bleibt es; wie laut der Inhalt insgesamt
% steht, ist eine Entscheidung pro ZSF. Zwei Faktoren wuerden die beiden Fragen
% vermischen (styles/README → Vereinigungstest).
\providecommand{\ZSFContentScale}{1}
\makeatletter
\newcommand{\ZSF@stepScale}[1]{%
  \setlength{\ZSF@fontSize}{#1\dimexpr\f@size pt\relax}%
  \setlength{\ZSF@fontSkip}{#1\dimexpr\f@baselineskip\relax}%
  \fontsize{\the\ZSF@fontSize}{\the\ZSF@fontSkip}\selectfont}
% Bei Faktor 1 gar nicht erst umschalten: Ein \fontsize auf denselben Wert
% kostet eine Fontauswahl und kann durch die pt-Rundung minimal danebenliegen.
\newcommand{\ZSF@applyContentScale}{%
  \ifdim\ZSFContentScale pt=1pt\else\ZSF@stepScale{\ZSFContentScale}\fi}
\newcommand{\ZSFfontContent}{\normalsize\ZSF@applyContentScale\ZSFApplyDisplaySkips}
\newcommand{\ZSFfontContentDense}{\footnotesize\ZSF@applyContentScale\ZSFApplyDisplaySkips}
\makeatother

% ── Freier Fliesstext: die dritte Inhaltsrolle ───────────────────────
% `runintext` ist der einzige Inhalts-Baustein OHNE eigene Schriftrolle — er
% erbt schlicht die Umgebungsgrösse. Damit war der Fliesstext das einzige
% Stück Satz, das man nur über einen rohen \fontsize erreichen konnte, und
% genau dort landete historisch der Notausgang in main.tex.
%
% Als eigene Rolle ist es eine Entscheidung statt eines Nebeneffekts: Wie laut
% steht die verbindende Prosa neben dem Bausteininhalt, der sie umgibt? In
% einer box-lastigen ZSF trägt die Box den prüfungsrelevanten Inhalt und der
% Fliesstext verbindet — dann darf er eine Spur leiser sein. In einer
% textlastigen ist es umgekehrt und der Faktor bleibt auf 1.
%
% Gerechnet wird von der KLASSENGRÖSSE aus, nicht vom aktuellen Font: So
% multiplizieren sich Grundgrösse und Rollenfaktor sauber (\ZSF@scaledFont
% legt \ZSFBodyScale selbst davor), statt sich zu verdoppeln, wenn die Rolle
% in einem bereits umgeschalteten Kontext aufgerufen wird.
\providecommand{\ZSFFlowTextScale}{1}
\makeatletter
\newcommand{\ZSFfontFlow}{%
  \ZSF@scaledFont{\ZSFFlowTextScale}{\ZSF@classFontSize pt}{\ZSF@classBaselineSkip}%
  \ZSFApplyDisplaySkips}
\makeatother

% ── Auszeichnungen INNERHALB der Bausteine ──
% Vier Rollen, die vorher als rohes \bfseries bzw. \small in den Bausteinen
% selbst standen. Als benannte Makros, weil sonst dieselbe Entscheidung an
% mehreren Stellen liegt: Wer die Tabellenköpfe entfetten will, müsste sonst
% 20_tables anfassen, wer die Listenmarken entfettet 60_boxes — und wer beides
% konsistent halten will, müsste es wissen (rules/10_architecture → Gestaltung
% an genau einem Ort).
%
% Kopfzelle einer ZSFtable (\ZSFhead).
\newcommand{\ZSFfontTableHead}{\bfseries}
% Marke einer ZSFlist und Beschriftung eines \ZSFsep-Blockwechsels — beide
% sind dieselbe Rolle: eine kurze fette Beschriftung innerhalb eines Bausteins.
\newcommand{\ZSFfontBlockLabel}{\bfseries}
% Text, der per \textVorBox/\textNachBox an eine Box gebunden wird.
\newcommand{\ZSFfontBoxBinding}{\small}
% Inline-Pille für Stolperfallen (\ZSFdanger).
\newcommand{\ZSFfontDangerTag}{\sffamily\bfseries}
% Zielnummer eines Verweises — \ZSFref, \ZSFsectionref und der Index-Locator
% (66_index). rules/55_index sagt zu, die Locators seien konsistent zu \ZSFref;
% mit drei rohen Kopien derselben Auszeichnung war das nur zufällig wahr.
\newcommand{\ZSFfontRefLocator}{\upshape\bfseries}

% Dezente Anmerkung — die EINE Quelle für Schrift und Farbe aller Anmerkungen
% in und an einer Box: \formulanote, \ZSFformulaline, \ZSFfigcaption und die
% Formelform von \textVorBox/\textNachBox. rules/20_boxes sagt zu, dass die
% Anmerkungs-Makros sich in der Platzierung unterscheiden und nicht im Stil;
% diese Zeile ist der Ort, an dem die Zusage eingelöst wird.
\newcommand{\ZSFBoxNote}[1]{{\ZSFfontNote\color{MutedTextColor}#1}}

% Beschriftung IN Diagrammen: TikZ-Knoten, Achsen-/Kurvenlabels und
% Annotationen an Formelteilen. Ersetzt lokales \tiny/\scriptsize an
% genau diesen Stellen — Diagramm-Inhalt ist Inhalt, seine Schriftgrösse
% ist Gestaltung und gehört deshalb hierher.
%
% ZWEI Stufen, weil ein Diagramm zwei Beschriftungsrollen hat und die Wahl
% pro Stelle fällt: die reguläre Marke (Achse, Kurve, Knoten) und die
% gedrängte Zusatzangabe (Mass an einer engen Kote, Index an einem kleinen
% Knoten), die neben der regulären stehen muss, ohne sie zu verdecken. Ohne
% die zweite Stufe entsteht genau dort das lokale \tiny, das dieser Block
% ersetzen soll.
%
% Für ein ganzes Bild einmal setzen statt pro Knoten:
%   \begin{tikzpicture}[font=\ZSFfontDiagramLabel] … \end{tikzpicture}
% Einzelne Knoten übersteuern das dann mit `font=\ZSFfontDiagramLabelSmall`.
%
% Die beiden Stufen sind ROLLEN, kein Massband: Sie sagen »reguläre Marke« und
% »gedrängte Zusatzangabe«, nicht 5.8bp und 4.6bp. Wie gross eine Rolle
% tatsächlich ausfällt, hängt am Fach — ein diagrammlastiges Blatt mit
% gedrängten Zeichnungen braucht beide Stufen kleiner, ohne dass die Rollen
% ihre Bedeutung verlieren. Deshalb steht davor ein Faktor und nicht eine
% dritte Stufe: Mehr Regler ist richtig, mehr Namen fast nie
% (rules/02_ai_mandate → Katalog-Ökonomie).
%
% \ZSFDiagramLabelScale wird wie die übrigen Faktoren in preamble.tex gesetzt,
% VOR diesem \input.
\providecommand{\ZSFDiagramLabelScale}{1}
% Gerechnet wird oben in \ZSF@scaledFont — dieselbe Stelle wie für jede andere
% Schriftrolle. Der Rollenfaktor \ZSFDiagramLabelScale multipliziert dort die
% Grundgrösse der ZSF: Ein diagrammlastiges Blatt dreht seine Beschriftungen
% kleiner, ohne dass sie aufhören, dem Satz der ZSF zu folgen.
\makeatletter
\newcommand{\ZSFfontDiagramLabel}{%
  \sffamily\ZSF@scaledFont{\ZSFDiagramLabelScale}{5.8bp}{6.4bp}}
\newcommand{\ZSFfontDiagramLabelSmall}{%
  \sffamily\ZSF@scaledFont{\ZSFDiagramLabelScale}{4.6bp}{5.1bp}}
\makeatother

% Semantische Inhaltsmakros (konsistente Skalierung und einfache Anpassung)
%
% \ZSFkeyword{...} — hebt zentrale Fachbegriffe im Fliesstext hervor.
% Verwendung:
%   * nur in \begin{runintext}, \begin{defbox}, \begin{warnbox} und
%     aehnlichen Text-Boxen, NIEMALS in Formeln, \formulanote oder
%     Überschriften (\SubsubsectionBar, \SubsectionBar, \StartChapter).
%   * sparsam einsetzen (Richtwert 1-3 pro runintext wegen der schmalen
%     Spalten) -- es sollen echte Fachbegriffe (Gesetze, Effekte,
%     Bauteile, zentrale Grössen) sein, keine generischen Substantive
%     oder Verben. Die Pro-Kapitel-Dichte ist bewusst KEIN fester Wert;
%     tests/check_index.sh warnt bei Über-Indexierung (Schwelle
%     ZSF_INDEX_DENSITY_MAX). Massgeblich ist rules/25_structure_markers.md.
% Rendering ist bewusst schlicht als \textbf{...}; Änderungen zentral
% hier vornehmen (z.B. Farbe, Underline) -- nicht in Kapiteln.
%
% Die Signatur ist INVARIANT gegenüber der Modulwahl: Stern und optionales
% Argument existieren immer, auch ohne 66_index. Sonst schluckt das Makro in
% einer ZSF ohne Index den Stern als Pflichtargument und \ZSFkeyword*{X}
% rendert als fettes "*" gefolgt vom Klartext "X" — ohne Fehler, ohne Warnung
% und ohne Linter-Treffer. Ein Baustein, dessen Aufrufform von einem Schalter
% abhängt, den der Kapitelautor nicht sieht, ist keine stabile API
% (rules/02_ai_mandate → gleiche Sache, gleiche Bedienung).
% 66_index ersetzt nur den Rumpf und ergänzt das Schreiben; die Form bleibt.
\NewDocumentCommand{\ZSFkeyword}{s o m}{\ZSFkeywordStyle{#3}}

% \ZSFlabel{...} — neutrale fette Beschriftung OHNE Fachbegriff-Semantik:
% Listenköpfe, Zeilenmarken, „Links & Ressourcen:". Existiert, weil sonst
% \ZSFkeyword* dafür zweckentfremdet wird — dessen Bedeutung ist aber
% „Fachbegriff, aber nicht indexieren", nicht „einfach fett". Ein Marker,
% der zwei Dinge heisst, verliert seine Aussage (rules/25_structure_markers).
\newcommand{\ZSFlabel}[1]{\ZSFlabelStyle{#1}}

% Die beiden Darstellungs-Hooks der Marker — getrennt, weil \ZSFkeyword und
% \ZSFlabel verschiedene Dinge bedeuten und eine ZSF den Unterschied ändern
% können soll, ohne ein einziges Kapitel anzufassen.
%
% BEIDE schlicht fett, ohne Farbe. \ZSFkeyword stand bis hierher fett in der
% KAPITELFARBE, mit dem Argument, der primäre Scan-Anker müsse sich von jeder
% anderen Fettung abheben. Das Argument ist am Satz gescheitert: In laufendem
% Text bringt eine eingefärbte Vokabel Unruhe in den Absatz, ohne dass die
% Farbe etwas sagt — sie wiederholt nur, in welchem Kapitel man ohnehin gerade
% ist. Wer den Absatz liest, wird gestört; wer ihn überfliegt, findet die
% Fettung auch ohne Farbe.
%
% Farbe bleibt damit im System das, was sie vorher schon war und wofür sie
% trägt: Kapitel-Identität auf Flächen (Balken, Akzente), Wegweiser zum ZIEL
% eines Verweises (\ZSFref) und semantische Zuordnung in Formeln (\ZSFmhl*).
% Als blosse Betonung im Fliesstext ist sie ausdrücklich nicht vorgesehen.
%
% Die Unterscheidung der beiden Marker liegt jetzt allein in der Bedeutung —
% \ZSFkeyword ist ein Fachbegriff und landet im Register, \ZSFlabel ist eine
% neutrale Beschriftung und tut das nicht. Das ist der Unterschied, der beim
% Schreiben zählt; er war nie der Farbe zu entnehmen.
\newcommand{\ZSFkeywordStyle}[1]{\textbf{#1}}
\newcommand{\ZSFlabelStyle}[1]{\textbf{#1}}

% \ZSFInk statt \textcolor: Der Skriptverweis steht genauso oft IM Titel eines
% Balkens wie im Boxinhalt (\SubsectionBar{Reihen\ZSFScriptRef{51}}). Ein hart
% gesetztes \textcolor behielte dort sein Grau und stuende damit lesbar-schwach
% auf der farbigen Flaeche — waehrend jeder andere farbtragende Inline-Marker
% die Kontrastfarbe erbt. \ZSFInk macht ihn zum Teil desselben Vertrags:
% im Boxinhalt gedaempft wie bisher, auf einer besitzenden Flaeche deren Tinte
% (rules/25_structure_markers -> Ink-Vertrag).
\newcommand{\ZSFScriptRef}[1]{\hfill\ZSFInk{\ScriptRefColor}{\ZSFfontBoxBinding (S.#1)}}

% \ZSFhl{...} — Hervorhebung des EINEN prüfungskritischsten Fakts pro Box
% (fett + unterstrichen). Bewusst sparsam: hoechstens einmal pro Block.
% Klare Rollentrennung:
%   * Fachbegriff/Scan-Anker        -> \ZSFkeyword
%   * Stolperfalle/kritische Ausnahme -> \ZSFdanger
%   * Folgerung                      -> \ZSFconclusion
%   * sonst die EINE Schlüsselaussage -> \ZSFhl
% \uline statt \underline: \underline steckt seinen Inhalt in eine
% unzerbrechliche Box — jede Hervorhebung, die länger als eine Zeile ist,
% sprengt damit die 50-mm-Spalte. \ZSFhl markiert aber laut Zweck eine ganze
% Aussage, nicht ein Wort. \uline bricht mit dem Absatz um.
% Fett als Deklaration, nicht als Argument von \uline: ein \textbf{...} als
% Argument macht den Inhalt für ulem wieder zu einer Einheit.
%
% Die Darstellung liegt wie bei \ZSFkeyword und \ZSFlabel in einem eigenen
% Hook. Ohne ihn erwischte »die Hervorhebung dieser ZSF umstellen« zwei von
% drei Markern — der dritte war der einzige mit fest verdrahteter Erscheinung,
% und man sah es erst am Satz (styles/README → Auszeichnungs-Hooks).
\newcommand{\ZSFhlStyle}[1]{{\bfseries\uline{#1}}}
\newcommand{\ZSFhl}[1]{\ZSFhlStyle{#1}}

% Papierorientierte Referenzrollen:
%   * \ZSFref{label}: hilfreicher Querverweis mit farbiger Abschnittsnummer.
%   * \ZSFsectionref{label}: kompakte Zielnummer für lokale Router/Tabellen.
% Ziel-Label via optionales Argument von \SubsubsectionBar/\SubsectionBar/\StartChapter setzen:
%   \SubsectionBar[sec:matrix-inverse]{Inverse}
% Farbe = Palette-Farbe des ZIELkapitels (Wegweiser-Prinzip): wird automatisch
% aus der Kapitelnummer der Referenz abgeleitet ("6.6" -> ChapterColor6).
% Bei Kapiteln jenseits der Palette wird derselbe Modulo-Lookup wie bei
% \StartChapter genutzt, damit gewrappte Kapitel und Verweise konsistent sind.
% Aufrecht + fett statt kursiv für maximale Lesbarkeit bei 6.7bp.
\makeatletter
\def\ZSF@refExtractFirst#1#2\@nil{#1}
\def\ZSF@refExtractChap#1.#2\@nil{#1}
% Bestimmt \ZSF@refTargetColor aus dem Label; Fallback auf die aktive
% Kapitelfarbe im ersten latexmk-Pass (Label noch undefiniert).
\newcommand{\ZSF@refComputeTargetColor}[1]{%
  \@ifundefined{r@#1}{%
    \edef\ZSF@refTargetColor{\chaptercolor}%
  }{%
    \edef\ZSF@refTargetNum{\expandafter\expandafter\expandafter
      \ZSF@refExtractFirst\csname r@#1\endcsname\@nil}%
    \edef\ZSF@refTargetChap{\expandafter\ZSF@refExtractChap\ZSF@refTargetNum.\@nil}%
    \ZSFSetPaletteColorNameFromNumber{\ZSF@refTargetChap}{\ZSF@refTargetColor}%
  }%
}
% Die Zielfarbe läuft über \ZSFInk. Ohne das steht ein Verweis auf das eigene
% Kapitel im Titel einer \SubsectionBar als ChapterColorN auf \chaptercolor!75 —
% derselbe Farbton bei voller Sättigung, also unlesbar, ohne dass Build oder
% Linter etwas melden. Auf einer Fläche mit eigener Kontrastfarbe erbt der
% Verweis diese und bleibt lesbar; der Farb-Wegweiser gilt dort, wo er auch
% ankommt — im Boxinhalt (40_colors_structure → Ink-Vertrag).
\newcommand{\ZSFref}[1]{%
  \begingroup
  \ZSF@refComputeTargetColor{#1}%
  \hyperref[#1]{{\ZSFfontNote\ZSFfontRefLocator
    \ZSFInk{\ZSF@refTargetColor}{($\rightarrow$\,\ref*{#1})}}}%
  \endgroup
}
\newcommand{\ZSFsectionref}[1]{%
  \begingroup
  \ZSF@refComputeTargetColor{#1}%
  \hyperref[#1]{{\ZSFfontRefLocator\ZSFInk{\ZSF@refTargetColor}{\ref*{#1}}}}%
  \endgroup
}
\makeatother

\newcommand{\ZSFconclusion}[1]{%
  \ZSFInterlude{\ZSFspaceS}%
  \noindent$\Rightarrow$ #1%
  \ZSFInterlude{\ZSFspaceS}%
}

